% TOP_PROBLEM: {"id": 20001055, "title": "Local edge budgets for Thomassé's nonneighbor conjecture", "category_id": 2, "category": {"id": 2, "name": "combinatorics", "display_name": "Combinatorics", "description": "Counting problems, graph theory, discrete structures.", "slug": "combinatorics", "order_index": 2, "created_at": "2026-07-31T15:26:25.671Z"}, "status": "partially_solved", "problem_number": "AIM-COMBINATORICS-0180", "set_id": 14, "statusQualification": "Uses the intended simple-digraph degree convention; unrestricted parallel arcs would make the literal degree statement false."}
% FIRST_POSTED: 2026-10-01
% ENTRY_KIND: statement_only
% UnsolvedMath ID 20001055; source catalogue code AIM-COMBINATORICS-0180
% !TeX program = lualatex
% Definition and sources only; this file is not an LLM solution attempt.
\documentclass[11pt,a4paper]{article}
\usepackage[margin=25mm]{geometry}
\usepackage{fontspec}
\setmainfont{lmroman10-regular.otf}[BoldFont=lmroman10-bold.otf,
 ItalicFont=lmroman10-italic.otf,BoldItalicFont=lmroman10-bolditalic.otf]
\usepackage{amsmath,amssymb,mathtools}
\usepackage{unicode-math}
\setmathfont{latinmodern-math.otf}
\AtBeginDocument{\let\setminus\smallsetminus}
\usepackage{xurl,hyperref}
\hypersetup{colorlinks=true,urlcolor=blue}
\setcounter{secnumdepth}{0}
\setlength{\parindent}{0pt}
\setlength{\parskip}{5pt}
\setlength{\emergencystretch}{3em}
\begin{document}
\section*{Local edge budgets for Thomassé's nonneighbor conjecture}\label{20001055}
{\small UnsolvedMath ID 20001055 · AIM-COMBINATORICS-0180 · Combinatorics · AIM Workshop Problem Lists}
\subsection{Definitions and mathematical statement}
Let $D=(V,E)$ be a nonempty finite simple digraph on $n$ vertices
having no directed cycle of length one, two, or three. In particular,
it has no loops or opposite arcs. For $v\in V$, define its non-neighbours
by
\[
 Z(v)=\{w\in V\setminus\{v\}:v\to w\notin E
                                      \text{ and }w\to v\notin E\}.
\]
Thus a non-neighbour has no adjacency to $v$ in either direction.
Let $d^+(v)$ and $d^-(v)$ count its outgoing and incoming neighbours.
The three sets of outgoing neighbours, incoming neighbours, and
non-neighbours partition $V\setminus\{v\}$.

Thomass\'e's conjecture is that some vertex satisfies
\[
                             d^+(v)\le |Z(v)|.
\]
Equivalently, there exists $v$ with
$2d^+(v)+d^-(v)\le n-1$.
The exclusion of directed triangles concerns cyclic orientations;
transitive triangles are allowed. The statement makes no bound on
the size of the digraph and does not assume that its degrees are
balanced or that it is strongly connected.
A sink automatically provides a witness, but the question applies
also to graphs with positive minimum outdegree. The non-neighbour
count excludes the vertex itself. The simplicity convention is
necessary when degree means number of arcs: repeated parallel arcs
could increase the outdegree without changing the non-neighbours or
creating a short cycle.
\subsection{Short English statement}
Does every directed graph without a cycle of length at most three have a vertex with no more outgoing neighbours than non-neighbours?
\subsection{Sources}
\begin{itemize}
\item[S1] ULAM.AI and the UnsolvedMath Contributors, supplied problems.json, record 20001055 (AIM-COMBINATORICS-0180), AIM Workshop Problem Lists. Catalogue identity and source transcription; CC BY 4.0.
\url{https://huggingface.co/datasets/ulamai/UnsolvedMath}.
\item[S2] American Institute of Mathematics, The Caccetta-Haggkvist conjecture, item 6. Original workshop question underlying AIM-COMBINATORICS-0180.
\url{https://aimath.org/WWN/caccetta/caccetta.pdf}.
\end{itemize}
\end{document}
